\documentclass[12pt,a4paper]{article}
\usepackage{amsmath}
\usepackage{geometry}
\usepackage{booktabs}
\usepackage{hyperref}
\geometry{margin=1in}

\title{\textbf{Reviewer Instructions}\\
\large Independent Evaluation Package\\
\normalsize Hubble Tension Manuscript — MNRAS Resubmission}
\author{Eric D. Martin \\ ORCID: 0009-0006-5944-1742}
\date{2026-04-01}

\begin{document}
\maketitle
\thispagestyle{empty}

\section*{Purpose of This Package}

This package is provided for independent evaluation by a reviewer
with no prior knowledge of this work. You should evaluate the
resubmission using \textbf{only the materials contained in this
package}. Do not consult external references, prior conversations,
or supplementary materials beyond what is included here.

\section*{Package Contents}

\begin{tabular}{ll}
\toprule
File & Contents \\
\midrule
\texttt{0\_instructions.pdf} & This document \\
\texttt{1\_original\_submission.pdf} & The manuscript as originally submitted to MNRAS \\
\texttt{2\_rejection.pdf} & The MNRAS editorial decision letter \\
\texttt{3\_our\_corrections.pdf} & Our analysis of each rejection issue and how it was addressed \\
\texttt{4\_resubmission.pdf} & The revised manuscript prepared for resubmission \\
\bottomrule
\end{tabular}

\section*{Your Task}

\subsection*{Step 1 — Read the original submission (Document 1)}
Read the original manuscript. Note any claims that are unclear,
unsupported, or that you cannot follow without external material.

\subsection*{Step 2 — Read the rejection letter (Document 2)}
The editor raised five specific issues. Read them carefully.
Note whether you agree or disagree with each issue based solely
on Document 1.

\subsection*{Step 3 — Read our corrections analysis (Document 3)}
Read our explanation of how each issue was addressed.
Do you agree that each issue has been adequately resolved?
Flag any issue you believe is still open.

\subsection*{Step 4 — Read the resubmission (Document 4)}
Read the revised manuscript. For each of the five editor issues,
confirm whether the resubmission addresses it. Be specific:
cite the section and paragraph that resolves each point.

\subsection*{Step 5 — Deliver your verdict}

Answer the following questions:

\begin{enumerate}
    \item Can you follow the core scientific argument from start to
    finish using only Document 4, without needing any external
    material (e.g.\ the full N/U Algebra framework, the UHA patent
    filing, prior publications)?

    \item Does the resubmission directly and adequately address
    each of the five editor objections? List any that remain open.

    \item Is there any claim in Document 4 that a professional
    cosmologist could reasonably reject on scientific grounds
    as stated? If so, identify it and explain why.

    \item Is the manuscript at a level of presentation and
    evidence suitable for a professional international
    astronomy journal?

    \item \textbf{Overall verdict:} Ready to resubmit, or does
    specific additional work remain?
\end{enumerate}

\section*{Constraint: No Outside Material}

The test is whether the manuscript stands on its own. If you find
yourself wanting to reference the UHA patent, the full N/U Algebra
framework, or any other external document to understand or defend
a claim, that is a failure of the manuscript --- not a failure of
this review process. Flag it explicitly.

\section*{Author's Position}

The author's position, stated plainly: the Hubble tension is not a
crisis in physics. It is a predictable artifact of comparing two
$H_0$ estimates computed within incompatible distance frameworks,
inflated by 5$\sigma$ significance that conflates a coordinate-frame
difference with a genuine physical discrepancy. The horizon-normalized
coordinate $\xi = d_c/d_H$ isolates the artifact from the real
residual. Applied to public data, it shows 93.3\% of the tension is
artifactual and 3\% is real. The 3\% residual is a matter density
discrepancy, not an expansion rate discrepancy.

That position must be evaluable, challengeable, and falsifiable
from Document 4 alone. If it is not, the manuscript is not ready.

\end{document}
